\documentclass[11pt]{article}
\usepackage[paperwidth=210mm,paperheight=297mm,margin=16mm]{geometry}
\usepackage{fontspec,amsmath,array,multirow,graphicx}
\usepackage{polyglossia}
\setmainlanguage{latin}\setotherlanguages{arabic,greek}
\setmainfont{Linux Libertine O}[Ligatures=TeX]
\newfontfamily\arabicfont{Amiri}[Script=Arabic]
\newfontfamily\greekfont{FreeSerif}[Script=Greek]
\newfontfamily\NallinoSigns{NallinoSigns.otf}[Path=./fonts/]
\newcommand{\AbjadZero}{{\NallinoSigns\symbol{"E000}}}
\newcommand{\AbjadOver}[1]{\leavevmode\vbox{\hrule height .45pt\kern1.2pt\hbox{#1}}}
\newcommand{\Name}[1]{{\addfontfeatures{LetterSpace=12}#1}}
\pagestyle{empty}\setlength{\parindent}{0pt}
\renewcommand{\arraystretch}{1.18}

\begin{document}
{\large \textarabic{٢٢٧}}\par\vspace{12mm}
\begin{Arabic}\begin{center}
{\large تنبيه}\\[8mm]
قال كَرْلو نالّينو المعتني بضَبْط هذا الكتاب وتصحيحه قد تمّت الابواب\\
كلّها فتليها في نسخة الأسْكوريال الجداول وامّا نحن فإنّما استخرجنا\\
منها ما يتعلَّق بالتأْريخ والجغرافيا واسماء الكواكب الثابتة واجرَيْنا\\
فيه الطبع بدون إصلاح ما وقع في حروف الجُمَّل من الخطأ\\
والتصحيف فواللّٰهِ ما اكثرَ هذه الأغْلاطَ . ومن يُرِدْ\\
تصحيحها فلْيُراجِع ترجَمتنا اللاتينيّة لهذا الكتاب\\
التي طبَعْنا فيها ايضًا بقيّة الجداول المُشْتَمِلة\\
على اعداد فقط\\[6mm]
وممّا يجِب تنبيه القارئ اليه أنّ حروف الجُمَّل مَعْناها في الجداول\\
كمعناها في النسخة الأسْكورياليّة يعني على مَذهب اهل\\
المغرِب فلذلك \AbjadOver{ص} عبارة عن ستّين و \AbjadOver{ض} عن\\
تسعين و \AbjadOver{س} عن ثلثمائة و \AbjadOver{ظ} عن ثمانمائة\\[6mm]
تم تم\\[2mm]
تم
\end{center}\end{Arabic}
\end{document}
